\chapter{Diseño de la solución}
\label{cap:diseno}

Este capítulo presenta las decisiones que definen la solución. El diseño busca separar las tareas que requieren interpretar lenguaje de aquellas que pueden resolverse mediante reglas exactas. Primero se describe la arquitectura general; después se justifican el contrato de datos, el esquema del \textit{IRS}, la red de agentes, la validación y la abstracción de proveedores. Los detalles de código se reservan para el capítulo 7.

\section{Arquitectura general del sistema}
\label{sec:arquitectura_general}

El sistema recibe una petición en lenguaje natural y produce uno de dos resultados: una \textit{RFQ} validada o un error que explica por qué no puede emitirse. Para cumplir este contrato, el flujo se divide en cinco etapas: ingesta, clasificación, extracción, validación y mapeo.

La ingesta conserva la petición original y añade una fecha de referencia para resolver expresiones relativas. El orquestador clasifica el producto como \texttt{IRS} o \texttt{UNSUPPORTED}. Si está soportado, el especialista extrae los siete campos obligatorios, deriva las convenciones permitidas y calcula los calendarios. El validador comprueba después la completitud y coherencia. Solo una operación válida llega al mapeador determinista, que genera la \textit{RFQ} final.

Cuando el validador detecta un error corregible de extracción, devuelve el diagnóstico al especialista. El bucle se limita a cinco pasadas. Si falta información que la petición nunca proporcionó, el proceso termina porque una nueva llamada no puede recuperarla. La clasificación queda fuera del bucle.

La arquitectura conserva un tercer agente para medir la fidelidad de serialización. Este agente se ejecuta después del mapeador, recibe los mismos términos validados y trata de reproducir su salida. Su respuesta no modifica la \textit{RFQ} emitida. Así se mide la capacidad del modelo sin convertirla en una dependencia operativa.

La configuración y la telemetría son transversales. Las instrucciones, el modelo y los recursos de cada agente se declaran fuera de la lógica financiera. Cada llamada registra proveedor, modelo, etapa, respuesta, consumo, latencia y coste. Esta separación permite comparar modelos sin cambiar el esquema ni las reglas.

La Figura~\ref{fig:arquitectura_general} resume el diseño.

\begin{figure}[H]
    \centering
    \fbox{
        \begin{minipage}{0.86\textwidth}
            \centering
            \textbf{Petición en lenguaje natural}\\[4pt]
            $\downarrow$\\[4pt]
            \textbf{Orquestador: clasificación}\\[2pt]
            \begin{tabular}{P{5.2cm} P{5.2cm}}
                $\swarrow$ \texttt{UNSUPPORTED} & \texttt{IRS} $\searrow$ \\
                Error: producto no soportado & Especialista de producto
            \end{tabular}\\[4pt]
            $\downarrow$\\[4pt]
            \textbf{Extracción, convenciones y calendarios}\\[4pt]
            $\downarrow$\\[4pt]
            \textbf{Validación determinista}\\
            {\small Error corregible: diagnóstico al especialista; máximo cinco pasadas}\\[4pt]
            $\downarrow$\\[4pt]
            \textbf{Mapeador determinista} $\rightarrow$ \textbf{RFQ emitida}\\[4pt]
            $\downarrow$\\[4pt]
            \textbf{Agente proto experimental} $\rightarrow$ \textbf{Métrica de fidelidad}\\[5pt]
            \rule{0.76\textwidth}{0.4pt}\\[3pt]
            {\small Telemetría: configuración, respuestas, tokens, latencia, coste y resultado}
        \end{minipage}
    }
    \caption[Arquitectura general del sistema]{Arquitectura general y posición experimental del agente proto}
    \label{fig:arquitectura_general}
    \floatfoot{Fuente: elaboración propia.}
\end{figure}

El criterio de diseño es sencillo: los modelos interpretan lenguaje; el código decide si los términos son admisibles y construye la salida. Esta frontera reduce el efecto probabilístico donde existen reglas deterministas y facilita atribuir cada fallo a una etapa.

\section{Justificación del uso de Protocol Buffers}
\label{sec:justificacion_protobuf}

La salida necesita un contrato que defina campos, tipos y relaciones entre las patas. Se consideraron JSON, FpML y \textit{Protocol Buffers}. JSON ofrece legibilidad y soporte amplio, pero necesita un esquema adicional para imponer la estructura. FpML proporciona una representación sectorial extensa, aunque su cobertura multiproducto supera las necesidades del prototipo \citep{fpml-home}. Protobuf permite definir mensajes tipados y generar clases para diferentes lenguajes desde un único archivo \citep{protobuf-overview}.

\begin{table}[H]
\centering
\small
\caption{Alternativas consideradas para el contrato de datos}
\label{tab:alternativas_contrato_datos}
\begin{tabular}{|P{2.7cm}|P{4.5cm}|P{5.6cm}|}
\hline
\textbf{Alternativa} & \textbf{Ventaja} & \textbf{Decisión} \\
\hline
JSON con esquema & Legible y ampliamente soportado. & Descartado porque obliga a mantener el formato y su contrato por separado. \\
\hline
FpML & Estándar específico y amplio para derivados. & Descartado por introducir estructuras fuera del alcance de un único producto. \\
\hline
Protocol Buffers & Tipos, enumeraciones, anidamiento y generación de código. & Elegido como contrato técnico reducido y compartido por los componentes. \\
\hline
\end{tabular}
\floatfoot{Fuente: elaboración propia a partir de \textcite{fpml-home} y \textcite{protobuf-overview}.}
\end{table}

La elección prioriza adecuación al alcance, no superioridad general. Protobuf representa de forma explícita los términos comunes y las dos patas, conserva la presencia de campos opcionales y permite analizar las respuestas de los agentes. El formato de texto facilita la inspección durante el experimento; una integración posterior podría utilizar la representación binaria sin cambiar el esquema.

El contrato estructural no sustituye la validación financiera. Puede exigir que el nocional sea numérico, pero no que sea positivo; puede almacenar fechas, pero no comprobar su orden; y no sabe si un índice es compatible con la divisa. Por ello, protobuf define la forma admisible y el validador determina si una instancia concreta es coherente.

\section{Separación entre esquema e instancia}
\label{sec:esquema_instancia}

El diseño distingue tres niveles. El esquema responde qué forma puede adoptar una \textit{RFQ}; la instancia contiene los valores de una petición; y la validación determina si esos valores representan una operación admisible.

\texttt{pricing.proto} es el contrato estable. Define el mensaje raíz, el \textit{IRS}, sus patas, tipos y enumeraciones. No contiene convenciones específicas de EUR o USD ni datos de una operación. Estas reglas permanecen fuera del esquema porque pueden evolucionar sin alterar la estructura básica.

Una instancia de cotización utiliza \texttt{PAR\_RATE\_QUOTE} y omite el tipo fijo. Una valoración utiliza \texttt{VALUATION} e incorpora el tipo pactado. Ambas comparten el mismo esquema; el validador comprueba que propósito y presencia del tipo sean coherentes.

La Figura~\ref{fig:esquema_instancia_validacion} muestra esta separación.

\begin{figure}[H]
    \centering
    \fbox{
        \begin{minipage}{0.84\textwidth}
            \centering
            \textbf{Esquema: \texttt{pricing.proto}}\\
            {\small Mensajes, campos, tipos y enumeraciones}\\[5pt]
            $\downarrow$\\[5pt]
            \textbf{Instancia: \texttt{RFQ.textproto}}\\
            {\small Valores extraídos para una petición concreta}\\[5pt]
            $\downarrow$\\[5pt]
            \textbf{Validación determinista}\\
            {\small Completitud, coherencia temporal y compatibilidad financiera}
        \end{minipage}
    }
    \caption[Separación entre esquema, instancia y validación]{Responsabilidades del contrato estructural, la operación concreta y las reglas financieras}
    \label{fig:esquema_instancia_validacion}
    \floatfoot{Fuente: elaboración propia.}
\end{figure}

Esta separación permite versionar el contrato, comparar instancias campo a campo y modificar convenciones sin redefinir el producto. También impide considerar válida una salida solo porque pueda analizarse como protobuf.

\section{Diseño del esquema de RFQ para IRS}
\label{sec:diseno_esquema_irs}

El esquema sigue la estructura de dos patas descrita en la Definición 2.19 de \textcite[p.~59]{ausin2025quantitative}. El mensaje raíz \texttt{RFQ} contiene un identificador y una instancia de \texttt{InterestRateSwap}. No se incluyó una estructura multiproducto porque la implementación evaluada solo trata \textit{IRS}; añadirla sin un segundo producto habría aumentado la complejidad sin aportar una capacidad medible.

\texttt{InterestRateSwap} agrupa el propósito, nocional, divisa, dirección, fechas principales y curva de descuento. También contiene \texttt{FixedLeg} y \texttt{FloatingLeg}. La curva de descuento permanece en el nivel común porque descuenta los flujos de ambas patas; la curva de estimación se sitúa dentro de la pata flotante porque solo proyecta sus fijaciones. Esta organización aplica el marco de doble curva descrito en la sección 2.4.6.2 de la misma referencia \citep[sec.~2.4.6.2]{ausin2025quantitative}.

La pata fija incluye el tipo cuando corresponde, la base, la frecuencia y las fechas de pago. La pata flotante incluye la naturaleza del tipo, índice, tenor cuando aplica, diferencial, base, frecuencia, curva de estimación y calendario. \texttt{FloatingRateType} distingue \texttt{IBOR} de \texttt{OVERNIGHT\_COMPOUNDED}; el segundo no utiliza tenor.

Los siete campos obligatorios son propósito, fechas de valoración, inicio y vencimiento, divisa, nocional y dirección. El tipo fijo no es obligatorio de forma general: debe faltar en una cotización y aparecer en una valoración. Los campos escalares se declaran como opcionales para distinguir ausencia y valor predeterminado.

Los calendarios se almacenan en dos vectores independientes. Esta decisión permite frecuencias distintas y periodos rotos. Cada elemento representa el final de un periodo y el último debe coincidir con el vencimiento. El esquema no incluye fechas de fijación ni calendarios de festivos, simplificaciones coherentes con el modelo de referencia utilizado en el trabajo \citep[p.~63]{ausin2025quantitative}.

La Figura~\ref{fig:jerarquia_esquema_irs} resume la jerarquía.

\begin{figure}[H]
    \centering
    \fbox{
        \begin{minipage}{0.9\textwidth}
            \centering
            \textbf{RFQ}\\
            {\small \texttt{rfq\_id}}\\[4pt]
            $\downarrow$\\[4pt]
            \textbf{InterestRateSwap}\\
            {\small propósito, nocional, divisa, dirección, fechas, curva de descuento}\\[5pt]
            \begin{tabular}{P{6.1cm} P{6.1cm}}
                $\swarrow$ \textbf{FixedLeg} & \textbf{FloatingLeg} $\searrow$ \\
                tipo fijo opcional & naturaleza del tipo e índice \\
                base y frecuencia & tenor opcional y diferencial \\
                fechas de pago & base, frecuencia y curva de estimación \\
                & fechas de pago
            \end{tabular}
        \end{minipage}
    }
    \caption[Jerarquía del esquema de IRS]{Jerarquía y ubicación de los campos principales del esquema}
    \label{fig:jerarquia_esquema_irs}
    \floatfoot{Fuente: elaboración propia.}
\end{figure}

El esquema prioriza una representación rigurosa del producto evaluado frente a una generalización no utilizada. Una futura ampliación deberá añadir cada producto con su propia estructura y reglas, sin forzar sus términos dentro del mensaje de \textit{IRS}.

\section{Diseño de la red de agentes}
\label{sec:diseno_red_agentes}

La red asigna una responsabilidad a cada componente. El orquestador clasifica; el especialista interpreta y estructura; el validador comprueba; el mapeador emite; y el agente proto mide la serialización. La Tabla~\ref{tab:responsabilidades_componentes} resume sus fronteras.

\begin{table}[H]
\centering
\small
\caption{Responsabilidades de los componentes del flujo}
\label{tab:responsabilidades_componentes}
\begin{tabular}{|P{3.0cm}|P{4.8cm}|P{5.0cm}|}
\hline
\textbf{Componente} & \textbf{Responsabilidad} & \textbf{No realiza} \\
\hline
Orquestador & Clasificar como \texttt{IRS} o \texttt{UNSUPPORTED}. & Extraer términos, validar o valorar. \\
\hline
Especialista & Extraer campos, derivar convenciones y calcular calendarios. & Decidir por sí solo la validez final. \\
\hline
Validador & Comprobar completitud y coherencia; decidir si procede reintentar. & Interpretar lenguaje libre o emitir la RFQ. \\
\hline
Mapeador & Construir la \textit{RFQ} desde términos validados. & Inferir o corregir términos financieros. \\
\hline
Agente proto & Reproducir experimentalmente la serialización. & Determinar la salida operativa. \\
\hline
\end{tabular}
\floatfoot{Fuente: elaboración propia.}
\end{table}

El orquestador recibe solo las reglas necesarias para delimitar el producto. El especialista añade conocimiento financiero y el esquema. Esta distribución evita cargar al clasificador con detalles que no necesita y permite cambiar las instrucciones de una etapa sin modificar las demás.

El bucle se sitúa entre extracción y validación. Esta posición permite devolver un diagnóstico específico sobre campos o convenciones, pero deja fuera la clasificación inicial. El límite de cinco pasadas acota coste y latencia. Una ausencia de información no se considera corregible; un error en términos existentes sí puede dar lugar a otro intento.

El mapeador se mantiene en el camino operativo porque la serialización no requiere interpretación. El agente proto permite comprobar empíricamente esta decisión. La arquitectura multiagente no se presenta como superior de forma universal; la comparación directa con una solución monolítica queda fuera de la evaluación y se propone como trabajo futuro.

\section{Capa de validación}
\label{sec:capa_validacion}

La validación se ejecuta después de analizar y normalizar la salida del especialista y antes de mapear la operación. Su función es impedir que una estructura formalmente correcta avance con términos ausentes o incompatibles.

La Tabla~\ref{tab:grupos_validacion} agrupa las principales reglas.

\begin{table}[H]
\centering
\small
\caption{Grupos de reglas de validación}
\label{tab:grupos_validacion}
\begin{tabular}{|P{3.4cm}|P{9.4cm}|}
\hline
\textbf{Grupo} & \textbf{Comprobaciones principales} \\
\hline
Completitud & Presencia de los siete campos obligatorios y de los atributos necesarios de cada pata. \\
\hline
Propósito & Tipo fijo ausente en \texttt{PAR\_RATE\_QUOTE} y presente en \texttt{VALUATION}. \\
\hline
Coherencia básica & Nocional positivo, valores plausibles y fechas válidas y ordenadas. \\
\hline
Convenciones & Bases, frecuencias, tipo flotante, tenor, curvas y compatibilidad entre divisa e índice. \\
\hline
Calendarios & Fechas crecientes, sin duplicados, dentro del plazo y con último pago en el vencimiento. \\
\hline
\end{tabular}
\floatfoot{Fuente: elaboración propia.}
\end{table}

Las convenciones derivadas se contrastan con las reglas de EUR y USD. Si la petición enuncia una alternativa compatible, esta prevalece. La detección de que un término fue expresado se realiza mediante un vocabulario cerrado, por lo que reconoce formas conocidas pero no comprende cualquier redacción equivalente. La compatibilidad esencial entre divisa e índice no puede anularse como convención alternativa.

El informe de validación contiene los errores, los campos ausentes y una indicación de si procede reintentar. Los datos ausentes producen un final explícito; los errores corregibles regresan al especialista con la petición original y el diagnóstico. El proceso termina al obtener una operación válida, detectar una ausencia no recuperable o alcanzar cinco pasadas.

Protobuf y el validador cumplen funciones complementarias. El primero asegura que la salida tenga una estructura reconocible; el segundo comprueba su significado financiero. Solo después de ambos controles se genera la \textit{RFQ}.

\section{Abstracción del proveedor de modelo}
\label{sec:abstraccion_proveedor}

OpenAI y Anthropic utilizan interfaces distintas para enviar instrucciones, devolver respuestas e informar del consumo. Una capa común oculta estas diferencias y entrega al resto del sistema el texto generado, los tokens de entrada y salida, el total y los tokens servidos desde caché cuando el proveedor los comunica.

El proveedor se deduce del identificador del modelo y la configuración permanece separada de las instrucciones. Así puede ejecutarse el mismo agente con modelos de ambos proveedores sin modificar el esquema, el validador ni el flujo. Las credenciales se obtienen de la configuración local y no se almacenan en el repositorio.

La telemetría distingue una respuesta incorrecta de un fallo externo. Si el proveedor devuelve contenido que incumple el contrato, la ejecución cuenta como resultado del modelo. Si no devuelve respuesta por red, autenticación o disponibilidad, el error se registra y queda fuera de los agregados de capacidad.

Las diferencias de uso también se normalizan. OpenAI incluye los tokens de caché dentro de la entrada; Anthropic los comunica por separado. Las tarifas de entrada, caché y salida se conservan en un archivo independiente para poder calcular y, si es necesario, reconstruir los costes desde los tokens archivados.

La temperatura se trata como una capacidad que el proveedor puede rechazar. Los modelos evaluados no permitieron fijarla: OpenAI devolvió un error al recibir el valor y el cliente de Anthropic lo rechazó antes de la llamada. El sistema reintentó sin el parámetro. Esta adaptación permitió mantener el flujo común, aunque impide apoyar la reproducibilidad en una temperatura igual a cero.

En conjunto, la abstracción hace comparables las ejecuciones sin ocultar las diferencias relevantes. La lógica financiera opera sobre un contrato estable y la telemetría conserva la información necesaria para interpretar costes, latencias y errores de cada proveedor.
